%! Diagonalizing a Matrix — Unit 6, Lesson 6.2 — Lesson Plan
\documentclass[10pt]{article}
\usepackage{linalg-article}
\usepackage{linalg-boxes}
\usepackage{graphicx}   % -article does NOT load graphicx; needed for thumbnails/screenshots
\graphicspath{{images/}}
\newcommand{\TallMath}[1]{$\displaystyle #1\rule[-1.4em]{0pt}{3.2em}$}
\newcommand{\xx}{\mathbf{x}}
\newcommand{\zero}{\mathbf{0}}
\newcommand{\UnitNumberName}{Unit 6: Eigenvalues and Eigenvectors \quad}
\newcommand{\LessonNumberName}{Lesson 6.2: Diagonalizing a Matrix}

\begin{document}
\begin{center}
    {\Large\bfseries \CourseName: \SchoolYear \\
    {\small\bfseries \UnitNumberName \LessonNumberName}}
\end{center}

\begin{tcolorbox}[colback=blush, colframe=burgundy, boxrule=0.9pt, arc=2mm,
    left=2mm, right=2mm, top=1.5mm, bottom=1.5mm]
    \textbf{Primary Objective:} Students can \emph{diagonalize} a $2\times2$ matrix. Given (or having
    found) the eigenvalues and eigenvectors, they line the eigenvectors into the columns of $X$ and the
    eigenvalues down the diagonal of $\Lambda$, explain why $AX=X\Lambda$, and rebuild
    $A=X\Lambda X^{-1}$. They can use the factorization to compute powers quickly via
    $A^k=X\Lambda^k X^{-1}$ (raising each eigenvalue to the power), interpret diagonalization as ``$A$ acts
    like the diagonal matrix $\Lambda$ in eigenvector coordinates,'' and decide when a matrix \emph{cannot}
    be diagonalized (too few independent eigenvectors --- e.g.\ a shear with a repeated eigenvalue).
\end{tcolorbox}

\renewcommand{\arraystretch}{1.4}
\begin{skillbox}[Priority Ideas \& Skills]{goldbox}
\begin{tabularx}{\linewidth}{@{}X|X@{}}
\textbf{Priority Ideas \& Skills} & \textbf{Key Understandings (the ``why'')} \\[2pt]
\hline
\vspace{-6pt}
\begin{itemize}[leftmargin=1.1em, nosep, after=\vspace{2pt}]
  \item Build $X$ (eigenvectors) and $\Lambda$ (eigenvalues); show $AX=X\Lambda$.
  \item Diagonalize: $A=X\Lambda X^{-1}$.
  \item Compute powers with $A^k=X\Lambda^k X^{-1}$.
  \item Decide whether a matrix is diagonalizable.
\end{itemize}
&
\vspace{-6pt}
\begin{itemize}[leftmargin=1.1em, nosep, after=\vspace{2pt}]
  \item $A$ scales each eigenvector, so $A\cdot X$ scales the columns $=X\Lambda$.
  \item In eigenvector coordinates, $A$ \emph{is} the diagonal $\Lambda$.
  \item The inner $X^{-1}X$'s cancel, so only the $\lambda$'s get raised to the power.
  \item Diagonalizable $\Leftrightarrow$ enough independent eigenvectors ($X$ invertible).
\end{itemize}
\end{tabularx}
\end{skillbox}

\begin{skillbox}[{Vocabulary, Concepts, \& Theorems}]{greenbox}
\renewcommand{\arraystretch}{1.3}
\begin{tabularx}{\linewidth}{@{}l X@{}}
\textbf{Eigenvector matrix $X$} & The matrix whose columns are the eigenvectors of $A$. \\
\textbf{Eigenvalue matrix $\Lambda$} & The diagonal matrix of eigenvalues, in the same order as $X$'s columns. \\
\textbf{$AX=X\Lambda$} & $A$ scales each column of $X$ by its eigenvalue --- the same as multiplying $X$ by $\Lambda$. \\
\textbf{Diagonalization} & $A=X\Lambda X^{-1}$ (equivalently $\Lambda=X^{-1}AX$), valid when $X$ is invertible. \\
\textbf{Power rule} & $A^k=X\Lambda^k X^{-1}$; $\Lambda^k$ just raises each eigenvalue to the $k$. \\
\textbf{Diagonalizable} & Has $n$ independent eigenvectors; a repeated eigenvalue may fail (e.g.\ a shear). \\
\end{tabularx}
\end{skillbox}

\begin{fixedskillbox}[Activate Prior Knowledge \& Spiral Review]{blush}
    The Warm-Up rehearses the three moves that \emph{assemble} today's factorization:
    \textbf{(1)} building a matrix from column vectors (eigenvectors $\to$ columns of $X$),
    \textbf{(2)} the $2\times2$ \emph{inverse} $\frac{1}{ad-bc}\begin{bsmallmatrix}d&-b\\-c&a\end{bsmallmatrix}$
    (Unit~2, for $X^{-1}$), and \textbf{(3)} right-multiplying by a \emph{diagonal} matrix (which scales the
    columns). Debrief by naming the plan aloud: ``today we combine these three moves into
    $A=X\Lambda X^{-1}$.''
\end{fixedskillbox}

\begin{skillbox}[Hook]{blush}
    Recall Lesson~6.1: for $A=\begin{bsmallmatrix}2&1\\1&2\end{bsmallmatrix}$ we \emph{found}
    $\lambda=3,1$ with eigenvectors $\begin{bsmallmatrix}1\\1\end{bsmallmatrix},\begin{bsmallmatrix}1\\-1\end{bsmallmatrix}$.
    Pose the payoff question: \textbf{``Suppose you needed $A^{10}$ --- would you really multiply $A$ by
    itself ten times?''} Along an eigenvector $A$ acts like a single number $\lambda$; the goal today is to
    make the \emph{whole} matrix that simple. Line the eigenvectors into $X$ and the eigenvalues into a
    diagonal $\Lambda$, and $A=X\Lambda X^{-1}$ --- then powers collapse to $A^k=X\Lambda^k X^{-1}$.
\end{skillbox}

\begin{skillbox}[Lesson]{blush}
\begin{multicols}{2}
\textbf{1. Line them up: $AX=X\Lambda$.} Eigenvectors into the columns of $X$, eigenvalues down $\Lambda$.
Multiplying $A\cdot X$ column by column scales each eigenvector, giving $AX=X\Lambda$
($\begin{bsmallmatrix}1&1\\1&-1\end{bsmallmatrix}$, $\Lambda=\mathrm{diag}(3,1)$).

\textbf{2. Solve for $A$.} If $X$ is invertible, right-multiply by $X^{-1}$: $A=X\Lambda X^{-1}$ (and
$\Lambda=X^{-1}AX$). Use the Unit~2 inverse formula for $X^{-1}$ ($\det X=-2$, halves), and verify
the product rebuilds $\begin{bsmallmatrix}2&1\\1&2\end{bsmallmatrix}$.

\columnbreak

\textbf{3. Payoff --- powers.} $A^2=X\Lambda X^{-1}X\Lambda X^{-1}=X\Lambda^2X^{-1}$; the $X^{-1}X$
cancels, so $A^k=X\Lambda^k X^{-1}$. Compute $A^2=\begin{bsmallmatrix}5&4\\4&5\end{bsmallmatrix}$ and check
by squaring. The pipeline picture: rewrite $\to$ scale $\to$ rebuild.

\textbf{4. When it fails + reach.} Need enough independent eigenvectors. The shear
$\begin{bsmallmatrix}2&1\\0&2\end{bsmallmatrix}$ (repeated $\lambda=2$, one eigenvector) is \emph{not}
diagonalizable. \textbf{Road ahead:} 6.3 symmetric (always works, $\perp$ eigenvectors), 6.4 an application.
\end{multicols}
\end{skillbox}

\begin{skillbox}[Active Monitoring]{redbox}
Circulate during the notes practice and Tier work. Watch for and correct:
\begin{itemize}[leftmargin=1.2em, nosep]
  \item column/eigenvalue \emph{order} mismatches between $X$ and $\Lambda$;
  \item multiplying $\Lambda$ on the wrong side (it must be on the \emph{right} to scale columns);
  \item fraction slips inverting $X$ (reinforce the Unit~2 formula; the triangular practice matrices avoid fractions);
  \item computing $A^k$ by brute force instead of raising the $\lambda$'s in $\Lambda^k$.
\end{itemize}
Cold-call prompts: ``Why is $AX=X\Lambda$?'' ``What is $X^{-1}$ here?'' ``What is $\Lambda^k$, and why do
the inner terms cancel?''
\end{skillbox}

\begin{skillbox}[Group Work \& Differentiation]{redbox}
\textbf{Diagonalizing a Matrix} activity (tiered; mirrors the notes).
\begin{multicols}{3}
\textbf{Tier R --- Remediate.} Given the eigen-data, build $X$ and $\Lambda$, invert $X$
(fraction-free, $\det X=1$), and \emph{assemble} $A=X\Lambda X^{-1}$.

\columnbreak
\textbf{Tier A --- Approaching.} The full job on $\begin{bsmallmatrix}4&-2\\1&1\end{bsmallmatrix}$
($\lambda=2,3$): find the eigen-data, diagonalize (integer $X^{-1}$), then compute $A^2=X\Lambda^2X^{-1}$
and check by squaring.

\columnbreak
\textbf{Tier E --- Extension.} The shear $\begin{bsmallmatrix}2&1\\0&2\end{bsmallmatrix}$: a repeated
eigenvalue with only \emph{one} eigenvector direction --- justify why it is \emph{not} diagonalizable.
\end{multicols}
\end{skillbox}

\begin{skillbox}[Individual Work \& Assessment]{redbox}
\textbf{Exit Ticket} (independent, no notes) on $A=\begin{bsmallmatrix}1&4\\0&3\end{bsmallmatrix}$: build
$X$ and $\Lambda$ (eigenvalues $1,3$; eigenvectors $\begin{bsmallmatrix}1\\0\end{bsmallmatrix},\begin{bsmallmatrix}2\\1\end{bsmallmatrix}$),
diagonalize $A=X\Lambda X^{-1}$ (fraction-free, $\det X=1$), and a \textbf{justification check} --- why is
$A^k=X\Lambda^k X^{-1}$ so much easier than repeated multiplication? Sort into ``got it / arithmetic slip /
concept slip'' to decide whether 6.3 opens with a quick diagonalization re-warm.
\end{skillbox}

\begin{skillbox}[Reinforcement \& Extension]{goldbox}
\textbf{Homework:} a full diagonalization (symmetric, with a negative eigenvalue); a \emph{power} via
$A^4=X\Lambda^4X^{-1}$; a diagonalizable-or-not decision (repeated vs.\ distinct eigenvalues); reading the
eigenvalues of $A^3$ and $A^{-1}$ off $\Lambda$; and a justification of the method (the $X^{-1}X$
cancellation; why $X$ must be invertible). \textbf{Extension:} $A^4$ of the notes matrix (one power
further). \textbf{Preview:} Lesson~6.3, \emph{Symmetric Positive Definite Matrices} --- the family that
\emph{always} diagonalizes, with perpendicular eigenvectors.
\end{skillbox}
\end{document}
